% Fichier généré par build_notebook_appendix.py — ne pas éditer à la main.
\medskip\noindent\textbf{\Large Naive Bayes multinomial en MapReduce sur Spark --- Démonstration}\par

Ce notebook tourne \textbf{de bout en bout sur les PETITES données (SMS Spam)}, \textbf{sans cluster}\par
(Spark en mode \texttt{local[*]}). Il montre :\par

1. le chargement et la vectorisation (bag-of-words) ;\par
2. l'entraînement \textbf{en RDD} (\texttt{map}/\texttt{reduceByKey}) ;\par
3. l'entraînement \textbf{en DataFrames} (\texttt{explode} + \texttt{groupBy}/\texttt{agg} + broadcast) ;\par
4. la vérification que les \textbf{deux versions donnent des prédictions identiques} ;\par
5. la \textbf{comparaison à \texttt{sklearn.naive\_bayes.MultinomialNB}}.\par

\textit{Prérequis : venv installé (\texttt{requirements.txt}) et un JDK 8/11/17. Le petit jeu SMS Spam}\par
\textit{est téléchargé automatiquement par la première cellule s'il est absent.}\par
\medskip\noindent\textbf{\large 0. Configuration : rendre \texttt{src/} importable et télécharger les données si besoin}\par

\noindent\textbf{Entrée [1]}
\begin{lstlisting}[language=Python]
import os, sys, subprocess

# Racine du projet = dossier parent de notebooks/
ROOT = os.path.abspath(os.path.join(os.getcwd(), ".." if os.path.basename(os.getcwd()) == "notebooks" else "."))
SRC = os.path.join(ROOT, "src")
sys.path.insert(0, SRC)

# Télécharger le petit jeu SMS Spam s'il n'existe pas encore.
sms_csv = os.path.join(ROOT, "data", "sms_spam.csv")
if not os.path.exists(sms_csv):
    print("SMS Spam absent -> téléchargement ...")
    subprocess.run([sys.executable, os.path.join(ROOT, "data", "download_sms_spam.py")], check=True)
print("Données :", sms_csv, "->", "OK" if os.path.exists(sms_csv) else "MANQUANT")
\end{lstlisting}


\noindent\textit{Sortie [1]}
\begin{lstlisting}[style=nboutput]
Données : /Users/abdrahamanecamara/Documents/Dauphine/Projet ML/data/sms_spam.csv -> OK
\end{lstlisting}

\medskip\noindent\textbf{\large 1. Chargement et préparation des données}\par

On charge les SMS (étiquettes \texttt{ham}/\texttt{spam}), on découpe en train/test (stratifié,\par
reproductible), puis on \textbf{tokenise} et on construit le \textbf{vocabulaire} sur l'entraînement.\par
Chaque document devient une liste d'indices de mots (bag-of-words).\par

\noindent\textbf{Entrée [2]}
\begin{lstlisting}[language=Python]
import nb_common as C

texts, labels = C.load_sms_spam(sms_csv)
print(f"{len(texts)} messages | spam={labels.count('spam')} ham={labels.count('ham')}")

# Découpage train/test reproductible (80/20, stratifié).
X_tr, X_te, y_tr, y_te = C.train_test_split_texts(texts, labels, test_size=0.2, seed=42)

# Tokenisation + vocabulaire + vectorisation en indices de mots (partagé par les 2 versions).
data = C.prepare(X_tr, y_tr, X_te, y_te)
print(f"Vocabulaire : {data.vocab_size} mots")
print(f"Train / Test : {len(data.train)} / {len(data.test)} documents")
print("Classes :", data.idx_to_label)

# Exemple : un document vectorisé (classe, premiers indices de mots)
c0, w0 = data.train[0]
print("Exemple doc -> classe", data.idx_to_label[c0], "| indices mots:", w0[:15], "...")
\end{lstlisting}


\noindent\textit{Sortie [2]}
\begin{lstlisting}[style=nboutput]
5574 messages | spam=747 ham=4827
Vocabulaire : 7716 mots
Train / Test : 4459 / 1115 documents
Classes : ['ham', 'spam']
Exemple doc -> classe ham | indices mots: [4917, 947, 4936, 6783, 7376, 6900, 3496, 3446, 3446] ...
\end{lstlisting}

\medskip\noindent\textbf{\large 2. Démarrage de Spark (local, sans cluster)}\par

\texttt{get\_spark} sélectionne automatiquement un JDK compatible (17) et expédie nos modules\par
Python aux workers. Le mode \texttt{local[*]} utilise tous les cœurs de la machine.\par

\noindent\textbf{Entrée [3]}
\begin{lstlisting}[language=Python]
spark = C.get_spark(app_name="demo-sms", master="local[*]")
spark.sparkContext.setLogLevel("ERROR")
sc = spark.sparkContext
print("Spark", spark.version, "| parallélisme =", sc.defaultParallelism, "cœurs")
\end{lstlisting}


\noindent\textit{Sortie [3]}
\begin{lstlisting}[style=nboutput]
26/07/08 21:23:35 WARN Utils: Your hostname, Macbook-Abdrahamane.local resolves to a loopback address: 127.0.0.1; using 192.168.0.32 instead (on interface en0)
26/07/08 21:23:35 WARN Utils: Set SPARK_LOCAL_IP if you need to bind to another address
Setting default log level to "WARN".
To adjust logging level use sc.setLogLevel(newLevel). For SparkR, use setLogLevel(newLevel).
26/07/08 21:23:35 WARN NativeCodeLoader: Unable to load native-hadoop library for your platform... using builtin-java classes where applicable
Spark 3.5.8 | parallélisme = 8 cœurs
\end{lstlisting}

\medskip\noindent\textbf{\large 3. Version RDD --- \texttt{map} / \texttt{reduceByKey}}\par

Le cœur MapReduce : pour chaque occurrence d'un mot \texttt{w} dans un document de classe \texttt{c},\par
on émet \texttt{((c, w), 1)} (\textbf{map}), puis \texttt{reduceByKey(add)} additionne les occurrences\par
(\textbf{reduce}) pour obtenir \texttt{count\_\{w,c\}}. Le modèle (log-probas) est ensuite construit\par
puis \textbf{broadcasté} pour la prédiction.\par

\noindent\textbf{Entrée [4]}
\begin{lstlisting}[language=Python]
import nb_rdd

model_rdd = nb_rdd.train_rdd(sc, data.train, data.vocab_size, data.idx_to_label)
acc_rdd = nb_rdd.evaluate_rdd(sc, model_rdd, data.test)
print(f"[RDD] accuracy test = {acc_rdd:.4f}")

# Aperçu du modèle : priors log P(c)
for c, lab in enumerate(model_rdd.idx_to_label):
    print(f"  log P({lab}) = {model_rdd.log_prior[c]:.4f}")
\end{lstlisting}


\noindent\textit{Sortie [4]}
\begin{lstlisting}[style=nboutput]
[RDD] accuracy test = 0.9848
  log P(ham) = -0.1440
  log P(spam) = -2.0091
\end{lstlisting}

\medskip\noindent\textbf{\large 4. Version DataFrames --- \texttt{explode} / \texttt{groupBy}-\texttt{agg}}\par

Même algorithme, exprimé en relationnel : \texttt{explode(words)} produit une ligne par mot\par
(\textbf{map}), \texttt{groupBy(label, word).count()} agrège (\textbf{reduce}). La prédiction utilise une\par
\textbf{UDF} s'appuyant sur le modèle \textbf{broadcasté}.\par

\noindent\textbf{Entrée [5]}
\begin{lstlisting}[language=Python]
import nb_dataframe

model_df = nb_dataframe.train_dataframe(spark, data.train, data.vocab_size, data.idx_to_label)
acc_df = nb_dataframe.evaluate_dataframe(spark, model_df, data.test)
print(f"[DataFrame] accuracy test = {acc_df:.4f}")
\end{lstlisting}


\noindent\textit{Sortie [5]}
\begin{lstlisting}[style=nboutput]
[DataFrame] accuracy test = 0.9848
\end{lstlisting}

\medskip\noindent\textbf{\large 5. Vérification : RDD et DataFrames donnent des prédictions identiques}\par

\noindent\textbf{Entrée [6]}
\begin{lstlisting}[language=Python]
pred_rdd = nb_rdd.predict_rdd(sc, model_rdd, data.test)
pred_df = nb_dataframe.predict_dataframe(spark, model_df, data.test)

print("Accuracy identique :", acc_rdd == acc_df, f"({acc_rdd:.4f})")
print("Prédictions identiques (ligne à ligne) :", pred_rdd == pred_df)
assert pred_rdd == pred_df, "Les deux versions doivent prédire exactement pareil !"
print("OK ✅")
\end{lstlisting}


\noindent\textit{Sortie [6]}
\begin{lstlisting}[style=nboutput]
Accuracy identique : True (0.9848)
Prédictions identiques (ligne à ligne) : True
OK ✅
\end{lstlisting}

\medskip\noindent\textbf{\large 6. Comparaison avec \texttt{sklearn.naive\_bayes.MultinomialNB}}\par

On entraîne sklearn avec \textbf{le même vocabulaire et la même tokenisation}, et le même\par
lissage (\texttt{alpha=1}). Notre implémentation reproduit MultinomialNB : les prédictions\par
doivent coïncider.\par

\noindent\textbf{Entrée [7]}
\begin{lstlisting}[language=Python]
from sklearn.feature_extraction.text import CountVectorizer
from sklearn.naive_bayes import MultinomialNB
from sklearn.metrics import accuracy_score

# Même vocabulaire (celui construit sur le train) + même tokenizer que nos versions Spark.
cv = CountVectorizer(tokenizer=C.tokenize, vocabulary=data.vocab,
                     token_pattern=None, lowercase=False)
Xtr = cv.transform(X_tr)
Xte = cv.transform(X_te)

y_tr_idx = [data.label_to_idx[l] for l in y_tr]
y_te_idx = [data.label_to_idx[l] for l in y_te]

clf = MultinomialNB(alpha=1.0)
clf.fit(Xtr, y_tr_idx)
pred_sk = list(clf.predict(Xte))

acc_sk = accuracy_score(y_te_idx, pred_sk)
print(f"[sklearn]   accuracy test = {acc_sk:.4f}")
print(f"[RDD/DF]    accuracy test = {acc_rdd:.4f}")
print("Prédictions Spark == sklearn :", [int(p) for p in pred_rdd] == [int(p) for p in pred_sk])
\end{lstlisting}


\noindent\textit{Sortie [7]}
\begin{lstlisting}[style=nboutput]
[sklearn]   accuracy test = 0.9848
[RDD/DF]    accuracy test = 0.9848
Prédictions Spark == sklearn : True
\end{lstlisting}

\medskip\noindent\textbf{\large 7. Prédiction sur de nouveaux messages}\par

On applique le modèle RDD (identique au DataFrame) à quelques messages inédits.\par

\noindent\textbf{Entrée [8]}
\begin{lstlisting}[language=Python]
nouveaux = [
    "Congratulations! You have WON a free prize, claim now!!!",
    "Hey, are we still on for lunch tomorrow?",
    "URGENT: your account needs verification, click this link to win cash",
    "Can you send me the slides before the meeting?",
]
# Vectorisation identique au pipeline d'entraînement.
new_indices = [C.doc_to_indices(C.tokenize(t), data.vocab) for t in nouveaux]
for txt, idxs in zip(nouveaux, new_indices):
    c = C.predict_indices(model_rdd, idxs)
    print(f"[{data.idx_to_label[c]:>4}]  {txt}")
\end{lstlisting}


\noindent\textit{Sortie [8]}
\begin{lstlisting}[style=nboutput]
[spam]  Congratulations! You have WON a free prize, claim now!!!
[ ham]  Hey, are we still on for lunch tomorrow?
[spam]  URGENT: your account needs verification, click this link to win cash
[ ham]  Can you send me the slides before the meeting?
\end{lstlisting}


\noindent\textbf{Entrée [9]}
\begin{lstlisting}[language=Python]
spark.stop()
print("SparkSession arrêtée.")
\end{lstlisting}


\noindent\textit{Sortie [9]}
\begin{lstlisting}[style=nboutput]
SparkSession arrêtée.
\end{lstlisting}

